% ============================================================================
% CHAPITRE 2 — Étude préalable
% Structure alignée sur la table des matières cible
% ============================================================================
\openchapter{Étude préalable}

\section{Introduction}

Avant de détailler la conception et la réalisation du projet (chapitre~3), ce
chapitre propose une étude préalable du sujet : pourquoi une plateforme dédiée
est nécessaire, à qui elle s'adresse, quelles approches alternatives existent
et pourquoi elles ne répondent pas au besoin, ainsi que le cadrage logique du
problème à résoudre.

\section{Pourquoi ce projet, à qui il s'adresse, approches existantes}

\subsection{Pourquoi ce projet ?}

Le fonctionnement d'une association repose sur deux flux qui doivent rester
étroitement liés : la cotisation annuelle que chaque adhérent verse selon les
résultats de l'assemblée générale, et les projets associatifs (forage de
puits, construction ou rénovation de mosquées, aide aux veuves et aux
orphelins, etc.) que l'association finance grâce à ces cotisations, aux dons
de ses propres membres ou de bienfaiteurs, ainsi qu'aux dons de l'État et
d'autres associations. Avant de lancer un projet, le bureau constitue un
comité chargé de le mener de A à Z, et chaque membre du comité doit pouvoir
saisir les opérations qui relèvent de son rôle : montant des dons reçus,
reçus de paiement, factures des fournisseurs. À la fin du projet, le comité
doit produire un rapport détaillant les entrées et les dépenses, accompagné
des factures, bons et copies de chèques ; le bureau, de son côté, doit
présenter un rapport annuel aux adhérents.

Tant que ce suivi repose sur des tableurs et des dossiers papier, il devient
rapidement fragile : rien n'empêche une saisie en double ou une perte de
justificatif, rien ne garantit qu'un membre du comité ne voie que ce qui
relève de son rôle, et un adhérent n'a aucun moyen simple de connaître sa
propre situation (cotisation payée, dons effectués, solde dû). Un système
qui centralise ces opérations, applique des droits différenciés selon le rôle
de chaque membre et conserve une traçabilité systématique des justificatifs
répond directement à cette limite.

\subsection{À qui s'adresse ce projet ?}

Le projet s'adresse à l'ensemble des parties prenantes d'une association
telle que décrite dans son cahier des charges : le \textbf{bureau de
l'association} (président, vice-président, trésorier, vice-trésorier,
secrétaire général, vice-secrétaire général et conseillers), qui pilote
l'association et supervise l'ensemble des projets et des cotisations ; les
\textbf{comités de projet}, constitués ponctuellement par le bureau pour
mener un projet donné, dont les membres saisissent les opérations financières
de ce projet selon leur rôle au sein du comité ; et les \textbf{adhérents}
(abonnés), qui financent l'association par leur cotisation et peuvent, selon
le cas, être aussi des donateurs ou des membres de comité. Les huit rôles
identifiés et leurs droits respectifs sont détaillés au chapitre~3
(section~III).

\subsection{Approches existantes}

Trois grandes familles d'outils permettent, en théorie, de gérer les
cotisations et les projets d'une association :

\begin{itemize}
  \item \textbf{Le suivi manuel (tableurs et papier)} — rapide à mettre en
    place, sans coût de développement, mais sans droits d'accès
    différenciés, sans traçabilité systématique des justificatifs, et sans
    consultation en libre-service pour l'adhérent : chaque demande de
    situation personnelle repasse par le bureau.

  \item \textbf{Un logiciel de comptabilité générique} — apporte une
    tenue de compte robuste, mais n'est pas structuré autour de la notion
    de \emph{comité de projet} : il ne distingue pas nativement les droits
    d'un trésorier de comité de ceux d'un simple membre, ne rattache pas un
    don ou une dépense à un projet et à son cycle de vie propre (comité
    constitué, financement réuni, projet actif, clôturé), et ne propose
    aucun espace personnel pour un adhérent. Son coût de licence est en
    outre disproportionné pour une association à but non lucratif.

  \item \textbf{Une plateforme de gestion associative sur mesure} —
    conçue autour des besoins réels exprimés dans le cahier des charges du
    stage : cotisations annuelles, projets avec comités dédiés, saisie des
    opérations selon le rôle de chaque membre, justificatifs numérisés
    (reçus, factures, photos de chèques ou d'ordres de virement), rapport de
    clôture par projet et historique personnel pour chaque adhérent.
\end{itemize}

C'est cette troisième approche qui a été retenue : elle seule permet de
modéliser fidèlement la structure à deux niveaux de l'association (rôle
global dans le bureau, rôle local dans un comité de projet) et de répondre
point par point aux exigences du cahier des charges, ce qu'aucun outil
générique ne permet sans détournement de son usage prévu.

\section{Cadrage logique}

\subsection{Problématique et enjeux}

L'enjeu du projet est double : offrir au bureau et aux comités de projet un
outil qui centralise et sécurise la gestion des cotisations et des
opérations financières, tout en restant suffisamment simple pour être adopté
par des utilisateurs non spécialistes de l'informatique.

\subsection{Arbre à problèmes et arbre à solutions}

\begin{figure}[H]
  \centering
  \begin{tikzpicture}[node distance=0.9cm and 0.5cm]
    \node[problembox] (p1) {Suivi des cotisations et des dons dispersé sur papier et tableurs};
    \node[problembox, right=of p1] (p2) {Aucune séparation des droits selon le rôle dans le comité de projet};
    \node[problembox, right=of p2] (p3) {Absence de traçabilité centralisée des justificatifs (reçus, factures, chèques)};
    \node[concbox, below=1.4cm of p2] (pc) {Gestion associative peu fiable et chronophage};
    \draw[bluearrow] (p1.south) -- (pc.north -| p1.south);
    \draw[bluearrow] (p2.south) -- (pc.north);
    \draw[bluearrow] (p3.south) -- (pc.north -| p3.south);
  \end{tikzpicture}
  \caption{Arbre à problèmes : causes principales d'une gestion associative peu fiable.}
  \label{fig:arbre-problemes}
\end{figure}

\begin{figure}[H]
  \centering
  \begin{tikzpicture}[node distance=0.9cm and 0.5cm]
    \node[problembox] (s1) {Plateforme centralisée de gestion des cotisations et des dons};
    \node[problembox, right=of s1] (s2) {Droits d'accès différenciés par rôle global et par rôle de comité};
    \node[problembox, right=of s2] (s3) {Traçabilité systématique des opérations avec justificatifs numérisés};
    \node[concbox, below=1.4cm of s2] (sc) {Système de gestion associative fiable et traçable};
    \draw[bluearrow] (s1.south) -- (sc.north -| s1.south);
    \draw[bluearrow] (s2.south) -- (sc.north);
    \draw[bluearrow] (s3.south) -- (sc.north -| s3.south);
  \end{tikzpicture}
  \caption{Arbre à solutions : réponses apportées aux causes identifiées.}
  \label{fig:arbre-solutions}
\end{figure}

\subsection{Objectif général du projet}

Concevoir et développer une plateforme web permettant au bureau de
l'association de gérer ses adhérents, ses cotisations annuelles et ses
projets associatifs ; aux comités constitués pour chaque projet de saisir les
opérations financières selon les droits de leurs membres ; et à chaque
adhérent de consulter en libre-service son propre historique (cotisation,
dons, solde dû) — le tout avec une traçabilité complète des opérations et de
leurs justificatifs.

\section{Conclusion}

Ce chapitre a posé le cadre logique du projet : pourquoi une plateforme
dédiée est nécessaire, à qui elle s'adresse, pourquoi les approches
existantes n'y répondent pas, et quel objectif général elle poursuit. Le
chapitre suivant détaille les missions et travaux réalisés : la modélisation
conceptuelle du système retenu et sa mise en \oe uvre informatique.